\documentclass[aps,prb,reprint,amsmath,amssymb,superscriptaddress,longbibliography,nofootinbib,floatfix]{revtex4-2}
\usepackage[T1]{fontenc}
\usepackage{lmodern}
\usepackage{graphicx}
\usepackage{bm}
\usepackage{microtype}
\usepackage[hidelinks]{hyperref}
\graphicspath{{../../results/novelty_20261008/four_cell_window_audit/}{../../results/novelty_20261008/figures/}{../../results/novelty_20261004/rank_figures/}{../../results/novelty_20261004/rank_controls_figure/}{../../results/novelty_20261004/evidence_figures/}{../figures/}}
\newcommand{\Dloc}{\Delta_{\mathrm{loc}}}
\newcommand{\Wlow}{W_{\mathrm{low}}}
\newcommand{\Cmin}{C_{\mathrm{min}}}
\begin{document}
\raggedbottom
\title{Thermal stability of retuned interference zeros in a Heisenberg pentamer chain}
\author{Author Name}
\email{corresponding.author@example.edu}
\affiliation{Department and Institution, City, Postal Code, Country}
\date{\today}
\begin{abstract}
Destructive interference can suppress a spin spectrum at zero temperature,
but thermal occupation introduces additional amplitudes that must be
canceled together. We determine when exchange retuning can maintain this
cancellation in a weakly coupled Heisenberg pentamer chain.
For a local doublet on a uniform reflection-symmetric periodic chain,
retuning cancels the entire first-order probe.
The weight is of order $\beta^4$ or smaller as the intercell exchange
$\beta$ tends to zero at fixed $T/J_{\rm eff}$.
A local quartet instead generates a composite vector whose bond-spin
channels require incompatible dipole corrections.
At a specified interference root, an exact commutator moment proves a
nonzero leading thermal response for every finite even ring with at
least four cells and every $T/J_{\rm eff}>0$.
The result survives independent first-order SU(2)-invariant exchange
controls within each cell.
The four-cell leading probe nevertheless annihilates the ground state;
its first thermally activated lines occur at $2J_{\rm eff}$ and
$7J_{\rm eff}$ with intensity ratio $9:8$.
Complete low-band spectra of twenty microscopic spins reproduce these
lines and resolve the crossover to a fourth-order ground-state background.
Finite momentum acceptance adds form-factor weight whose variance survives
window-optimized retuning.
The local spin representation thus determines whether a tunable
ground-state interference zero can suppress the entire thermal response
within the weak-coupling control space.
\end{abstract}
\maketitle
\section{Introduction}

Destructive interference can make a collective spin excitation invisible
to a scattering probe. In cluster magnets, amplitudes from sites within
a unit cell combine into a momentum-dependent form factor.
Its zeros suppress intensity even when the corresponding excitation
remains in the spectrum~\cite{waldmann2003interference,haraldsen2011pentamer,tabrizi2021selection,tabrizi2021anisotropic}.
Weakly coupled clusters provide a useful setting for studying this
separation between excitation energies and their visibility.
Their local ground multiplets form effective spin chains, while
intracell interference organizes the folded spectral
weight~\cite{cheng2022trimer,li2026tetramer}.

Near a form-factor zero, corrections to the scattering operator become
essential. Such corrections modify dimer structure
factors~\cite{mueller2003structure} and introduce tensor contributions
when cluster multiplets mix~\cite{waldmann2005manyspin}.
Symmetry constrains these amplitudes but can leave their relative
strengths dependent on the Hamiltonian.
For isotropic clusters, repeated occurrences of a transition-operator
irrep permit independent reduced matrix elements and their
interference~\cite{tabrizi2021selection}.
This distinction also underlies the separation between systematic
extinctions and parameter-dependent zeros~\cite{tabrizi2026extinctions}.
An exchange parameter that cancels one amplitude therefore need not
cancel every transition involving the same symmetry labels.

Thermal occupation adds a further constraint on cancellation.
Excited initial states contribute additional
transitions, as described by effective-observable corrections and
thermal spectral-weight transfer in spin
ladders~\cite{exius2010vertex,honecker2016thermal}.
Thermal lifting of accidental zeros also occurs in magnetovibrational
interference~\cite{raymond2018interference}.
Here we ask whether exchange retuning can compensate the additional
amplitudes and keep the entire low-energy thermal response dark.
The answer requires both the corrected probe and the space of
amplitudes that physical controls can adjust.

We address this question in a Heisenberg pentamer chain whose local
ground multiplet can be a doublet or a quartet.
The doublet permits cancellation of the entire first-order probe on a
uniform reflection-symmetric periodic chain.
At a specified quartet root, a composite vector instead produces
bond-spin amplitudes that remain incompatible with simultaneous
thermal cancellation.
We establish this obstruction for every finite even ring with at least
four cells, including independent first-order SU(2)-invariant controls
inside each cell.
A four-cell ring makes the physical consequence explicit: its leading
ground-state response can vanish, while thermal occupation restores
two calculable groups of spectral lines.
Microscopic spectra test these predictions, and momentum-window
averaging identifies the background that accompanies finite acceptance.

\section{Model and low-energy observable}

Each cell contains five spin-$1/2$ sites at unit spacing. With $J_1=1$,
$\alpha=J_2/J_1$, and $\beta=J_3/J_1>0$, the Hamiltonian is
\begin{align}
 H_0&=\sum_{n=0}^{N-1}\big[
 \alpha(\bm S_{n,1}\!\cdot\!\bm S_{n,2}
       +\bm S_{n,4}\!\cdot\!\bm S_{n,5})\notag\\
 &\hspace{29mm}+\bm S_{n,2}\!\cdot\!\bm S_{n,3}
       +\bm S_{n,3}\!\cdot\!\bm S_{n,4}\big],\\
 H&=H_0+\beta V,\qquad
 V=\sum_{\langle n,n+1\rangle}\bm S_{n,5}\!\cdot\!\bm S_{n+1,1}.
 \label{eq:rank-ham}
\end{align}
The bond sum is specified as open or periodic in each calculation, and the
microscopic length is $L=5N$. We assume that the isolated cell has a single
ground multiplet of spin $s$, separated from all other multiplets by
$\Dloc>0$. The numerical examples below have $s=1/2$ or $3/2$.

Rotational symmetry within this irrep fixes the local projection,
\begin{equation}
 P_n\bm S_{n,a}P_n=c_a\bm T_n,\qquad
 c_a=c_{6-a},\qquad \sum_a c_a=1.
 \label{eq:rank-projection}
\end{equation}
Here $\bm T_n$ is a spin-$s$ operator. The leading Hamiltonian is an
antiferromagnetic effective chain with exchange $\beta c_1^2$.
The scattering intensity additionally depends on how the microscopic
probe projects into this effective spin space.

It is convenient to center the intra-cell Fourier phase:
\begin{align}
 B_n^z(q)&=\sum_{a=1}^5 e^{iq(a-3)}S_{n,a}^z,\\
 O_q&=\frac{1}{\sqrt{5N}}\sum_n e^{iQn}B_n^z(q),
 \quad Q=5q\pmod{2\pi}.
\end{align}
Changing the coordinate origin adds only a global phase to $O_q$ and does not
change a spectral weight. Equation~\eqref{eq:rank-projection} gives
\begin{align}
 P O_qP&=A(q,\alpha)D_Q,\notag\\
 A(q,\alpha)&=\sum_a c_a\cos[q(a-3)],\\
 D_Q&=\frac{1}{\sqrt{5N}}\sum_n e^{iQn}T_n^z.
 \label{eq:rank-amplitude}
\end{align}
The normalization of $D_Q$ is fixed by the microscopic probe and is
kept identical across all comparisons.

At zero temperature, we use the integrated inelastic low-band weight
\begin{equation}
 \Wlow=\frac{1}{n_G}\sum_{r=1}^{n_G}
 \sum_{0<E_m-E_G<\Dloc/2}|\langle m|O_q|G_r\rangle|^2.
 \label{eq:rank-weight}
\end{equation}
All transitions within the ground eigenspace are excluded.
The gap $\Dloc$ in this cutoff and in $g=\beta/\Dloc$ is evaluated at the
unshifted local root. For each finite system, the cutoff contains the
complete retained low band in the weak-coupling regime used below.
The microscopic expansion is taken at fixed $N$; band isolation is
checked separately for each numerical calculation.

\section{Exchange retuning and the effective probe}

\subsection{Scalar and quadrupolar local factors}

At a form-factor zero, virtual excitations outside the local ground
multiplet determine the leading scattering amplitude. Let $P=\prod_n P_n$
and $R=(1-P)/(H_0-NE_5)$, where $E_5$ is the isolated ground-multiplet
energy. At $A(q,\alpha_*)=0$, the effective probe becomes
\begin{align}
 (O_q)_{\mathrm{eff}}&=\beta K_q+O(\beta^2),\notag\\
 K_q&=-P O_q R V P-P V R O_q P.
 \label{eq:rank-probe}
\end{align}
The two terms describe a virtual excitation created by the probe or by
the intercell exchange. They are evaluated separately for non-Hermitian $O_q$.

A link can contribute only through an intermediate excitation on the cell
acted on by the probe: the local probe cannot de-excite both cells of a link.
For a Hermitian component of $B^z$, the relevant local tensor is
\begin{equation}
 M^{z\mu}=P_n(B^z R_n S_{\mathrm{end}}^\mu
                   +S_{\mathrm{end}}^\mu R_n B^z)P_n.
 \label{eq:rank-local-tensor}
\end{equation}
It is Hermitian and even under time reversal. Coupling two vectors permits
ranks zero, one, and two, but a Hermitian rank-one tensor within a single
spin irrep is odd under time reversal. Only the scalar and quadrupolar parts
remain. The real and imaginary Hermitian components of $B(q)$ can be treated
separately and recombined.

Define
\begin{equation}
 Q_n^{z\mu}=\tfrac12\{T_n^z,T_n^\mu\}
        -\tfrac13 s(s+1)\delta_{z\mu}I.
\end{equation}
On one orientation of a link the correction consequently has the form
\begin{equation}
 \delta B_n^z=a\,T_{n+1}^z
       +b\sum_\mu Q_n^{z\mu}T_{n+1}^\mu,
 \label{eq:rank-link}
\end{equation}
with a corresponding term for the opposite endpoint.
The second term couples a local quadrupole to a neighboring spin,
forming a composite vector under simultaneous spin rotations.

\subsection{Cancellation for a local doublet}

In a single spin-$1/2$ irrep, $Q_n^{z\mu}=0$ identically. Equivalently every
Hermitian local matrix is a linear combination of $I$ and $\bm T$; time
reversal leaves only $I$ in Eq.~\eqref{eq:rank-local-tensor}.
The entire first-order correction is thus a sum of neighboring one-spin
operators. On a uniform periodic chain at a commensurate momentum, translation
reduces this sum to $K_q=\gamma(q)D_Q$. Cell-centered reflection and real
exchange give $a_R=a_L^*$ for the two link orientations. Consequently
$\gamma=a_L e^{-iQ}+a_R e^{iQ}=2\operatorname{Re}(a_L e^{-iQ})$
is real in the same phase convention as $A$.

Suppose the local zero is simple, with
$A'=\partial_\alpha A(q,\alpha_*)\ne0$. A physical exchange retuning
$\alpha=\alpha_*+\kappa\beta$ gives
\begin{equation}
 (O_q)_{\mathrm{eff}}
   =\beta[\gamma+\kappa A']D_Q+O(\beta^2).
\end{equation}
Choosing $\kappa=-\gamma/A'$ cancels the first-order operator throughout
the retained Hilbert space. Every low-band transition therefore satisfies
\begin{equation}
 \Wlow=O(\beta^4)\quad\text{or smaller}.
 \label{eq:rank-doublet}
\end{equation}
This result applies to a single isolated doublet with SU(2) and
time-reversal symmetry, bilinear links, and the stated periodic and
reflection symmetries. Higher-order retuning can further reduce the
remaining weight; open endpoints are treated in
Appendix~\ref{app:ground-controls}.

\subsection{Residual for a quartet}

For $s=3/2$, the composite vector supplies an additional scattering
amplitude. To determine whether it survives in the ground-state response,
define $d=D_Q|G\rangle$ and
$k=K_q|G\rangle$, with the entire ground subspace projected out.
For degenerate ground states these vectors denote a direct sum, with norm
averaged over the ensemble used in Eq.~\eqref{eq:rank-weight}.
The leading weight after a real retuning is
\begin{equation}
 \Wlow=\beta^2\|k+\kappa A'd\|^2+O(\beta^3).
\end{equation}
Elementary least squares gives
\begin{align}
 \kappa_{\mathrm{opt}}
   &=-\frac{\operatorname{Re}\langle d,k\rangle}{A'\|d\|^2},\\
 \Cmin&=\|k\|^2-
       \frac{[\operatorname{Re}\langle d,k\rangle]^2}{\|d\|^2}.
 \label{eq:rank-minimum}
\end{align}
The minimum is the part of the source orthogonal to the accessible
dipole correction. It depends on the many-body state as well as on the
local tensor structure. The overlaps in the examples below are real to
numerical precision.

At $q=\pi$ the quartet branch provides a simple analytic example:
\begin{align}
 \alpha_*&=-\frac{2+\sqrt7}{6},\notag\\
 c_1=c_5&=\frac{25+4\sqrt7}{108},\quad
 c_2=c_4=\frac14,\notag\\
 c_3&=\frac{1-2\sqrt7}{27}.
 \label{eq:rank-exact-root}
\end{align}
These coefficients obey $c_1-c_2+c_3-c_4+c_5=0$.
Here $\Dloc=0.16622155$, $c_1^2=0.10855198$, and
$A'=0.10026088$. The local resolvent gives $a=-0.12123577$ and
$b=0.12361176$, with exact expressions and a proof that $b>0$ in
Appendix~\ref{app:exact-local-check}. Thus this root realizes a nonzero
composite contribution within the microscopic pentamer model.

\subsection{A bond-spin interpretation}

The composite vector also has a direct interpretation in terms of neighboring
total-spin channels. For two spin-$s$ sites, put
$h_{ij}=\bm T_i\cdot\bm T_j$, $D_{ij}=T_i^z-T_j^z$, and
$X_{ij}=\sum_\mu(Q_i^{z\mu}T_j^\mu-T_i^\mu Q_j^{z\mu})$.
Expanding the quadrupoles gives the exact identity
\begin{equation}
 X_{ij}=\tfrac12\{h_{ij},D_{ij}\}+\tfrac13s(s+1)D_{ij}.
 \label{eq:bond-identity}
\end{equation}
Within bond total spin $J$, $h_{ij}$ has eigenvalue
$[J(J+1)-2s(s+1)]/2$. The exchange-odd vector $D_{ij}$ connects only
adjacent bond-spin sectors. If $\Pi_J$ denotes their projector, then
\begin{equation}
 \Pi_J X_{ij}\Pi_{J-1}=
 \left[\frac{J^2}{2}-\frac{2s(s+1)}{3}\right]
 \Pi_J D_{ij}\Pi_{J-1}.
 \label{eq:bond-channel}
\end{equation}
For a doublet there is only $J=1\leftrightarrow0$, and this coefficient
vanishes. For a quartet, the $1\leftrightarrow0$, $2\leftrightarrow1$,
and $3\leftrightarrow2$ coefficients are $-2,-1/2,2$.
The link probe $-aD_{ij}+bX_{ij}$ therefore renormalizes these channels
differently, whereas exchange retuning adds the same multiple of $D_{ij}$.
One adjustment cannot cancel all three link-operator coefficients.
The ground state can nevertheless sample a restricted combination of
these channels and remain dark. Thermal occupation brings additional
initial states into play, making simultaneous cancellation a property
of the full low-energy dynamics.

\section{Thermal stability of a retuned zero}
\label{sec:thermal}

\subsection{An operator zero and a ground-state zero}

Temperature tests whether exchange retuning suppresses all low-energy
transitions or only those accessible from the ground state.
We keep $\tau=T/J_{\rm eff}$ fixed as
$\beta\to0$, where $J_{\rm eff}=\beta c_1(\alpha_*)^2$ and $k_B=1$.
Thus the thermal ensemble remains inside the local ground multiplets.
Put $h=\sum_i\bm T_i\cdot\bm T_{i+1}$ and let $\epsilon_i$ be its
dimensionless eigenvalues. For the Hermitian leading $Q=\pi$ probe,
write $D=D_\pi$, $K=K_q$, and define
\begin{align}
 R_\kappa&=K+\kappa A'D,\notag\\
 C_\kappa(\tau)&=\sum_{\epsilon_f>\epsilon_i}
 p_i\left|\langle f|R_\kappa|i\rangle\right|^2,
 \qquad p_i=\frac{e^{-\epsilon_i/\tau}}{Z}.
 \label{eq:thermal-weight}
\end{align}
This coefficient describes transitions between distinct effective energy
levels, with frequencies on the $J_{\rm eff}$ scale.
At fixed finite size, choose $\eta>0$ smaller than
the smallest nonzero level spacing of $h$, and let $W_{+,\eta}$ include
only low-band transitions with $\omega>\eta J_{\rm eff}$. Then
$W_{+,\eta}=\beta^2C_\kappa+O(\beta^3)$.
We call this the regular positive-frequency response.
Higher-order splittings within an equal-energy block instead produce
quasielastic lines at $\omega=O(\beta^2)$, outside this leading coefficient.
The minimum $C_{\min}(\tau)$ allows a separate real retuning at each
temperature, including any temperature-dependent displacement of the zero.

At any $\tau>0$ every state of a finite effective chain has positive
Boltzmann weight. Nonnegativity of the summands and Hermiticity imply
\begin{equation}
 C_\kappa(\tau)=0\quad\Longleftrightarrow\quad[h,R_\kappa]=0.
 \label{eq:thermal-criterion}
\end{equation}
Thus complete thermal cancellation requires the retuned probe to act
only within equal-energy blocks. Equation~\eqref{eq:thermal-criterion}
reduces the thermal question to an operator calculation.
For a doublet, the stronger identity
$K=\gamma D$ already proved above cancels $R_\kappa$ throughout the
retained space. Consequently $W_+=O(\beta^4)$ or smaller at every fixed
$\tau$ under the same periodic and reflection assumptions.

\subsection{A representation-dependent obstruction}

For the quartet root at $q=\pi$, define the unnormalized alternating spin
$\widetilde D=\sum_i(-1)^iT_i^z$ and
$Y=\sum_i(-1)^iX_{i,i+1}$, using Eq.~\eqref{eq:bond-identity}. Then
\begin{equation}
 R_\kappa=\frac{bY+(-2a+\kappa A')\widetilde D}{\sqrt{5N}}.
 \label{eq:thermal-source}
\end{equation}
Both $\widetilde D$ and $Y$ are spin vectors with momentum $\pi$,
odd under time reversal and even under the site-permutation reflection
$i\mapsto-i$. For $Y$, reflection reverses the bond orientation and
shifts its alternating sign, leaving their product unchanged.
Their different bond-channel amplitudes, rather than different symmetry
labels, control the cancellation. To test whether one dipole coefficient
can remove every inelastic matrix element, we evaluate the squared
commutator using a spectral moment~\cite{mueller1982sumrules}.
Write $u=-2a+\kappa A'$ in Eq.~\eqref{eq:thermal-source}.
For every finite even quartet chain with $N\ge4$, the uniformly traced,
two-sided second frequency moment is
\begin{align}
 M_2(u)&=\frac{1}{4^N}\operatorname{Tr}[h,R]^\dagger[h,R]\notag\\
 &=\frac{197}{16}b^2+\frac52(u+b)^2.
 \label{eq:quartet-moment}
\end{align}
The positive first term survives every choice of $u$.
Together with Eq.~\eqref{eq:thermal-criterion}, it proves
$C_{\min}(\tau)>0$ at every $\tau>0$ for these finite chains.
Even a ground state that can be made dark therefore acquires an
irreducible leading thermal response.

The proof uses orthogonality of local operator supports, so the same
minimum applies at every finite even $N\ge4$.
Appendix~\ref{app:general-moment} derives the spin-$s$ expression and
shows how its minimum vanishes for a doublet.
Here the uniform trace covers the retained effective space, corresponding
to $J_{\rm eff}\ll T\ll\Dloc$ in the fixed-size weak-coupling limit.
Combining its second and fourth moments also bounds a finite-frequency
part of the quartet weight away from zero, uniformly in even $N$
(Appendix~\ref{app:weight-bound}). This latter result takes the effective
uniform-trace limit before increasing the chain length.

\subsection{The accessible local-control space}

Independent exchange controls within each cell leave this obstruction
intact. At the quartet root $q=\pi$, any bounded
SU(2)-invariant perturbation of order $\beta$ confined to an individual
cell changes its projected vector probe only by a multiple of $T_i^z$.
The isolated cell contains a single spin irrep; its projected scalar
perturbation shifts all components equally. The leading uniform $h$ and
the coefficient $b$ are therefore unchanged. Even independent local
controls in every cell can only enlarge the leading probe to
\begin{equation}
 R_{\{u_i\}}=\frac{bY+\sum_i(-1)^iu_iT_i^z}{\sqrt{\ell N}},
 \qquad u_i\in\mathbb R.
 \label{eq:local-control-probe}
\end{equation}
Here $\ell=5$ fixes the normalization inherited from the microscopic probe.
These controls preserve the local irrep and the leading intercell
Hamiltonian; their changes to either $h$ or $b$ enter at higher order.

Translation symmetry separates the $Q=\pi$ component, whose dipole
coefficient is $\bar u=N^{-1}\sum_i u_i$, from all other operator momenta.
Their thermal spectral weights add without interference. For either the
positive-frequency weight or any fixed positive-frequency window,
\begin{equation}
 F(R_{\{u_i\}})=F(R_{\bar u})+
 F(R_{\{u_i\}}-R_{\bar u})\ge F(R_{\bar u}).
 \label{eq:local-control-projection}
\end{equation}
Here $F$ uses a Gibbs state of the uniform leading $h$, or its uniform
trace. The minimum over all real $u_i$ therefore equals the uniform
minimum. Spatial variation adds spectral weight at other momenta but
cannot compensate the residual $Q=\pi$ response.
Appendix~\ref{app:local-controls} gives the full proof.

\section{Thermally activated response and microscopic spectra}

\subsection{Exact activation from a dark four-cell ground state}

The four-cell quartet ring makes the contrast explicit: the retuned
leading probe annihilates its ground state, yet thermally occupied
states scatter. Write $\bm A=\bm T_0+\bm T_2$ and
$\bm B=\bm T_1+\bm T_3$, so that $h=\bm A\cdot\bm B$.
Its ground state has $(j_A,j_B,J)=(3,3,0)$ and energy $-12$;
the first excited multiplet is $(3,3,1)$ at $-11$.
The zero-temperature optimum is $\lambda_0=-2a-3b$, with
$\lambda=-\kappa A'$. The resulting source
$b(Y+3\widetilde D)/\sqrt{20}$ annihilates the ground state exactly,
although it does not commute with $h$.

The first thermally occupied triplet opens two groups of transitions.
Their final sectors are $(3,3,2)$ at $\omega/J_{\rm eff}=2$ and
$(3,1,2),(1,3,2)$ at $\omega/J_{\rm eff}=7$.
Summed over the three initial magnetic components, the squared matrix
elements divided by $b^2$ are respectively $27/50$, $6/25$, and $6/25$.
Consequently,
\begin{equation}
 C_{\min}(\tau)=\frac{51}{50}b^2e^{-1/\tau}
                  +o(e^{-1/\tau}).
 \label{eq:thermal-activation}
\end{equation}
The two groups have leading intensity ratio $9:8$. Reoptimizing the
exchange at each temperature changes only higher activation orders.
These four-cell lines identify the transitions that restore scattering
when the first excited triplet becomes populated.

\begin{figure*}[t]
\centering
\includegraphics[width=\textwidth]{thermal_evidence.pdf}
\caption{Thermal response at retuned interference zeros.
(a) Minimum leading positive-frequency coefficient in four- and six-cell
effective periodic chains. The doublet zero follows from cancellation of
the first-order operator. (b) Full microscopic two-cell weights
at $T/J_{\rm eff}=1.6$, normalized separately to their values at
$g=\beta/\Dloc=0.005$. Quartet and doublet slopes are $2.02549$ and
$4.00189$.
(c) Weights averaged over the uniform interval $q_0\pm\delta q$, relative
to their ideal central-momentum values, at $g=0.005$ and
$T/J_{\rm eff}=1.6$. The quartet and doublet central momenta are
$\pi$ and $3\pi/5$, respectively. Exchange parameters remain on their
central-momentum retuning trajectories. Both initial and final microscopic
states lie below $E_0+\Dloc/2$, while the partition function includes
the full spectrum. Lines connect computed points without smoothing.}
\label{fig:thermal-evidence}
\end{figure*}

\subsection{Microscopic spectra and the low-temperature crossover}

Full thermal traces on four- and six-cell effective chains give the
contrast in Fig.~\ref{fig:thermal-evidence}(a). At $\tau=1.6$, the quartet
minima are $0.00991078$ and $0.00968212$, respectively. For six cells the
latter is about 22 times its zero-temperature value, even after
temperature-dependent retuning.

The microscopic coupling dependence tests the predicted distinction
between second- and fourth-order weights.
We diagonalized all 1024 states of two coupled pentamers, including both
closing links and their corresponding effective exchange bonds.
The quartet root is Eq.~\eqref{eq:rank-exact-root}; the doublet uses
$q=3\pi/5$, $\alpha_*=1.09755143$, so both have $Q=\pi$.
Over $g=0.005,0.01,0.02$, the weights at $\tau=1.6$ have logarithmic
slopes $2.02549$ and $4.00189$ for the quartet and doublet, respectively
[Fig.~\ref{fig:thermal-evidence}(b)]. The quartet coefficient approaches
the effective prediction to within $1.18\%$ at the weakest coupling.
Appendix~\ref{app:exact-local-check} reports the temperature dependence
and finite-coupling errors.

The twenty-spin periodic model resolves the predicted four-cell activation.
We reconstructed its complete 256-state low band at
$g=0.000625,0.00125,0.0025$, keeping the exchange on the leading
zero-temperature trajectory $\kappa_0=1.28029750$.
At the weakest coupling the first
thermal gap is $1.000318J_{\rm eff}$; the two line-group centroids are
$2.000592J_{\rm eff}$ and $7.000309J_{\rm eff}$. Their unweighted
transition strengths, summed over the initial triplet and divided by
$\beta^2b^2$, are $0.536037$ and $0.482098$, within $0.74\%$ of
$27/50$ and $12/25$. Both approach their predictions as $g$ decreases
[Fig.~\ref{fig:four-cell-activation}(b)].

At finite coupling, a fourth-order ground-state background competes with
the exponentially small thermal contribution. At $\tau=0.1$, the
weakest-coupling weight exceeds the leading prediction by $25.9\%$;
ground-state transitions account for $21.0\%$ of the microscopic total.
This background scales as $\beta^{4.0015}$ over the three couplings and
is included in Fig.~\ref{fig:four-cell-activation}(a).
The response therefore crosses from thermal activation to a
ground-state background as temperature decreases at fixed $\beta$.

Higher-order splittings also produce quasielastic weight, amounting to
$1.68\%$ of the total at $\tau=1.6$.
The regular response with $\eta=0.05$ agrees with the effective
prediction within $0.23\%$ at the weakest coupling for
$\tau=0.2,0.4,1.6$. The resolved and unwindowed weights, together with
cutoff and solver checks, are specified in
Appendix~\ref{app:exact-local-check}.

\begin{figure*}[t]
\centering
\includegraphics[width=\textwidth]{four_cell_activation.pdf}
\caption{Direct four-cell microscopic check and low-temperature background.
(a) Resolved thermal weight for twenty physical spins with periodic
boundaries, $q=\pi$, fixed zero-temperature retuning, and
$\omega>0.05J_{\rm eff}$. Both initial and final states are in the
complete 256-state low band. The finite-coupling curves retain the
ground-state background. The effective curve is the full leading thermal
trace at the same retuning; the activation asymptote is
$(51/50)b^2\exp(-J_{\rm eff}/T)$.
(b) Strengths from the first thermally occupied triplet to the two predicted
final-state groups, divided by $\beta^2b^2(27/50)$ and
$\beta^2b^2(12/25)$, respectively. These strengths are summed over initial
magnetic components before thermal weighting. Lines connect calculated
points without smoothing.}
\label{fig:four-cell-activation}
\end{figure*}

\subsection{Momentum acceptance and window retuning}

A finite momentum interval samples nonzero form factors on either side
of the interference root, adding weight to the central response.
We averaged microscopic transition probabilities over a uniform window
around each root, initially keeping the central retuning fixed. At
$g=0.005,\tau=1.6$, a half-width $0.001\pi$ raises the quartet weight
by $0.094\%$, but the doublet weight by a factor of approximately
$2.93\times10^6$ [Fig.~\ref{fig:thermal-evidence}(c)].
The local doublet form factor has $\partial_qA=1.69998$ at its root,
whereas the quartet root at $q=\pi$ is stationary in momentum.
A fixed window therefore adds an ordinary form-factor background that
can overwhelm the doublet's very small fourth-order central signal.
At this fixed width, the doublet weight becomes nearly independent of
coupling, with logarithmic slope $0.00475$.

Optimizing the exchange for the entire window reduces, but does not
eliminate, this background. At the stationary quartet root, write
$x=q-\pi$, $A_{qq}=\partial_q^2 A(\alpha_*,\pi)$, and let $B$ be the
positive-frequency thermal weight of $D$. If $\kappa_0$ minimizes the
central response, the remaining source is thermally orthogonal to $D$.
The leading window-dependent dipole amplitude is
$\beta A'(\kappa-\kappa_0)+A_{qq}x^2/2$.
Retuning removes its mean, leaving the variance of $x^2$.
For a uniform half-width $\delta q=\rho\sqrt{\beta}$, at fixed finite
size and scaled temperature, this gives
\begin{align}
 \kappa_{\rm win}&=\kappa_0-\frac{A_{qq}\rho^2}{6A'}+o(1),\notag\\
 \frac{\overline W^{\rm opt}_{+,\eta}}{\beta^2}
 &\longrightarrow C_{\min}+\frac{A_{qq}^2B}{45}\rho^4.
 \label{eq:window-retuning}
\end{align}
Keeping $\kappa_0$ fixed gives $1/20$ in place of $1/45$ in the
background coefficient. Window retuning therefore leaves $4/9$ of this
leading background, set by the variance of the form factor.
In the ten-spin ring at $\tau=1.6$, the optimized microscopic
coefficients for $\rho=0.25,0.5,1$ agree with this limit within
$0.30\%$ at $g=0.00125$; deviations increase with coupling.

The chosen doublet root has nonzero $A_q=\partial_q A$, so its leading
window background is instead $A_q^2B(\delta q)^2/3$. Keeping it at
the fourth-order scale requires $\delta q=O(\beta^2)$, whereas the
stationary quartet root permits $O(\sqrt{\beta})$ at its quadratic
scale. The different tolerances follow from the slope and curvature
of these particular momentum zeros.
In the two-cell doublet, further exchange optimization cancels the
single low-band singlet--triplet amplitude to numerical precision.
Its weight for $\delta q=\rho\beta^2$ then approaches
$\beta^4 A_q^2B\rho^2/3$, within $0.055\%$ at the weakest tested coupling.
Thus the fourth power can also arise from momentum acceptance, even
when the central response has been suppressed further.

\section{Discussion and conclusion}

Exchange retuning distinguishes two physically different origins of a
dark low-energy spectrum. For a local doublet, the first-order correction
has the same operator form as the projected dipole probe.
A single retuning cancels it throughout the retained space, suppressing
both ground-state and thermal weights to fourth order or smaller.
For the pentamer quartet, virtual excitations generate a composite vector
with different amplitudes in different bond-spin channels.
The positive minimum of its commutator moment forces a leading thermal
response at every positive scaled temperature on finite even rings.

The four-cell ring shows why a zero-temperature spectrum alone cannot
distinguish these mechanisms. Its leading quartet probe annihilates the
ground state after retuning, although the full operator remains dynamical.
Population of the first triplet restores lines at $2J_{\rm eff}$ and
$7J_{\rm eff}$ with leading intensity ratio $9:8$.
The twenty-spin microscopic spectrum approaches these line positions
and strengths in the weak-coupling limit.
Thus thermal scattering tests the compatibility of cancellation across
initial states, beyond the transition amplitudes sampled by the ground state.

The control-space result gives this distinction a direct physical meaning.
Independent first-order SU(2)-invariant changes within each cell adjust
local dipole coefficients while preserving the leading uniform chain.
Their nonuniform components add weight at other operator momenta, leaving
the optimal $Q=\pi$ residual unchanged.
The obstruction is therefore tied to the accessible operator structure
under these local controls. Altering the leading intercell interactions,
introducing anisotropy, or changing the local multiplet would change
that structure and poses a different control problem.

At finite coupling and momentum acceptance, the thermal signal competes
with two calculable backgrounds. A fourth-order ground-state contribution
dominates at sufficiently low temperature, while averaging around the
interference root restores ordinary form-factor weight.
Optimizing exchange over the momentum window removes its mean shift but
leaves its variance. These effects set the scales needed to interpret
a suppressed spectrum and separate thermal activation from residual
ground-state scattering.

The results establish a weak-coupling distinction at fixed finite size
between cancellation of a probe and cancellation for one initial state.
The effective uniform-trace bound further retains finite-frequency
quartet weight uniformly in even chain length, with the trace limit
taken first. Extending that bound to fixed scaled temperature in the
thermodynamic limit would connect the finite-chain mechanism to bulk
thermal scattering. Within the regime established here, the local
spin representation and accessible exchange controls determine whether
a retuned interference zero survives thermal occupation.

\begin{acknowledgments}
[Author, affiliation, funding, and acknowledgment details are retained for
completion by the authors.]
\end{acknowledgments}
\section*{Data and code availability}
Calculation scripts, numerical tables, and computational protocols are
available in the accompanying project files. Public repository details
will be supplied by the authors.
\appendix
\section{Exact second moment for the dipole--composite probe}
\label{app:general-moment}

The commutator moment tests simultaneous cancellation by summing the
squared inelastic amplitudes with their squared frequencies.
Consider a uniform spin-$s$ nearest-neighbor periodic chain with even
$N\ge4$ and $s\ge1/2$. Let
$R=(bY+u\widetilde D)/\sqrt{\ell N}$, with the operators defined in
Sec.~\ref{sec:thermal} and normalization $\ell=5$ for the pentamer.
Set $c=s(s+1)$ and $d=2s+1$. We derive
\begin{align}
 M_2(u)&=\frac{1}{d^N}\operatorname{Tr}[h,R]^\dagger[h,R]\notag\\
 &=M_{2,\min}+\frac{(d^2-1)^2}{18\ell}
       \left[u+\frac{b(d^2-4)}{12}\right]^2,\\
 M_{2,\min}&=\frac{b^2c^2(4c-3)(86c-27)}{810\ell}.
 \label{eq:general-moment}
\end{align}
For $s=1/2$ the quadrupole vanishes, giving a zero minimum.
For $s\ge1$ and nonzero $b$, the minimum is positive.
The result applies to this dipole--composite probe; its coefficient $b$
must be determined from the microscopic cluster.

The derivation separates same-bond and adjacent-bond commutators.
Their different operator supports make their contributions orthogonal
in the uniform trace.
On one bond, the quadrupole--spin vector obeys
\begin{equation}
 X_{ij}=\tfrac12\{h_{ij},D_{ij}\}+\tfrac{c}{3}D_{ij},
 \qquad D_{ij}=T_i^z-T_j^z.
\end{equation}
In the $J-1\to J$ channel its coefficient relative to $D_{ij}$ is
$r_J=J^2/2-2c/3$. The squared $D^z$ matrix elements summed over magnetic
components are $J(d^2-J^2)/3$, and the corresponding energy difference is
$J$. Including both transition directions, define
$w_J=2J^3(d^2-J^2)/3$ for $1\le J\le d-1$. Thus the Gram entries of
$C_D=[h_{ij},D_{ij}]$ and $C_X=[h_{ij},X_{ij}]$ are sums of
$w_J$, $w_Jr_J$, and $w_Jr_J^2$:
\begin{align}
 G_{DD}&=\frac{d^2(d^2-1)^2}{18},\notag\\
 \frac{G_{DX}}{G_{DD}}&=\frac{d^2-4}{12},\qquad
 \frac{G_{XX}}{G_{DD}}=\frac{7(d^2-4)^2}{360}.
 \label{eq:general-bond-gram}
\end{align}
Evaluation of these finite power sums gives the minimized bond norm
$d^2(d^2-1)^2(d^2-4)^2/1440$.
For $s=3/2$, the Gram matrix is
\begin{equation}
 \begin{pmatrix}G_{DD}&G_{DX}\\G_{DX}&G_{XX}\end{pmatrix}
 =\begin{pmatrix}200&200\\200&560\end{pmatrix},
 \label{eq:thermal-gram}
\end{equation}
with minimized bond norm $360$.

The adjacent-bond contribution on sites $0,1,2$ is
\begin{equation}
 Z_{012}=[h_{12},X_{01}]-[h_{01},X_{12}].
\end{equation}
With repeated Cartesian indices summed, its explicit decomposition is
\begin{align}
 Z_{012}/i={}&\epsilon_{abc}
  (Q_0^{zb}T_1^cT_2^a+T_0^aT_1^cQ_2^{zb})\notag\\
 &-T_0^a(\epsilon_{bzc}Q_1^{ca}
                  +\epsilon_{azc}Q_1^{cb})T_2^b.
 \label{eq:three-site-commutator}
\end{align}
The three terms have local tensor-rank patterns $QTT$, $TTQ$, and
$TQT$, respectively. They are mutually orthogonal because
$\operatorname{tr}(T^aQ^{bc})=0$. Rotational invariance and the traces of
$(T^z)^2$ and $(T^z)^4$ give
\begin{align}
 \operatorname{tr}(T^aT^b)&=t\delta_{ab},\quad t=dc/3,\notag\\
 \operatorname{tr}(Q^{ab}Q^{ef})&=v(\delta_{ae}\delta_{bf}
 +\delta_{af}\delta_{be}-\tfrac23\delta_{ab}\delta_{ef}),\notag\\
 v&=dc(4c-3)/60.
\end{align}
Contracting the Levi-Civita symbols in
Eq.~\eqref{eq:three-site-commutator} gives $20/3$, $20/3$, and $20$
times $vt^2$ for those three patterns. Therefore
\begin{equation}
 \|Z_{012}\|_{\rm HS}^2=\frac{100}{3}vt^2
                      =\frac{5}{81}d^3c^3(4c-3).
 \label{eq:three-site-norm}
\end{equation}

All one-site partial traces of $C_D,C_X$ and all two-site partial traces
of $Z$ vanish. Distinct two-site and three-site supports are consequently
orthogonal in the full uniform trace. For a periodic chain with $N\ge4$,
there are $N$ distinct supports of each type. The alternating signs give
\begin{equation}
 M_2(u)=\frac{\|bC_X+uC_D\|_{\rm HS}^2}{\ell d^2}
       +\frac{b^2\|Z\|_{\rm HS}^2}{\ell d^3}.
\end{equation}
Completing the square with Eqs.~\eqref{eq:general-bond-gram} and
\eqref{eq:three-site-norm} yields Eq.~\eqref{eq:general-moment}.
The adjacent-bond contribution is independent of $u$.
Keeping only the same-bond contribution gives the quartet bound
\begin{equation}
 \min_u\frac{\|[h,R]\|_{\rm HS}^2}{4^N}
 \ge\frac{360b^2}{5\times16}=\frac92b^2,
 \label{eq:thermal-bound}
\end{equation}
which the exact result sharpens to $(197/16)b^2$.

The local identities and partial traces were checked with
exact rational matrices for $s=1/2,1,3/2,2,5/2,3$. Separate orthonormal
Cartesian-spin Kronecker matrices on the full $N=4,6$ chains verify the
whole quadratic for $s=1/2,1,3/2$ to floating-point accuracy.

The trace in Eq.~\eqref{eq:general-moment} includes both frequency signs;
the positive-only moment is half as large. At the four-cell ground-state
retuning $u=3b$, it gives $M_2=(837/16)b^2=0.79932812$.
The twenty-spin low-band transition sum, divided by
$\beta^2J_{\rm eff}^2$ with uniform initial weights, yields $0.79967945$
at $g=0.000625$, a $0.044\%$ deviation.
The deviations increase to $0.088\%$ and $0.177\%$ at the next two
couplings. The frequency-squared factor suppresses quasielastic
contributions in this leading limit, allowing the complete transition
sum to be used without a frequency cutoff.

\subsection{A size-uniform positive-frequency weight bound}
\label{app:weight-bound}

Use the normalized Hilbert--Schmidt norm on each operator's support,
$\|O\|_2^2=d^{-|\operatorname{supp}O|}\operatorname{Tr}O^\dagger O$.
The full uniform-trace spectral measure of $R$ is symmetric in frequency.
Define its moments $M_{2k}$ including both frequency signs; $M_0$ also
includes equal-energy matrix elements. Local tensor orthogonality gives
\begin{equation}
 M_0=\frac{1}{\ell}\left[\frac{c}{3}u^2
          +\frac{c^2(4c-3)}{27}b^2\right].
 \label{eq:uniform-M0}
\end{equation}
For $s\ge1$, comparison with Eq.~\eqref{eq:general-moment} yields
$M_2\ge(3c/2)M_0$. To see positivity without optimizing an irrational
eigenvalue, rescale $b,u$ by the square roots of their coefficients in
$\ell M_0$. The matrix for $\ell[M_2-(3c/2)M_0]$ in these coordinates has
diagonal entries $(61c-42)/30$ and $7c/6$, and determinant
$c(107c-54)/180$. All are positive for $c\ge2$.

For the fourth moment write
\begin{align}
 [h,R]&=\frac{1}{\sqrt{\ell N}}\sum_i(-1)^i B_i,\notag\\
 B_i&=bC_{X,i}+uC_{D,i}+bZ_{i,i+1,i+2}.
\end{align}
The two-site and three-site terms are orthogonal as above, giving
$\|B_i\|_2^2=\ell M_2$. Only the four bonds meeting sites $i,i+1,i+2$
contribute to $G_i=[h,B_i]$. Since $\|h_{ij}\|_{\rm op}=c$,
$\|G_i\|_2\le8c\|B_i\|_2$. Each traceless $G_i$ is supported within
the five-site interval from $i-1$ to $i+3$. Two disjoint intervals have
zero uniform-trace inner product. At most nine translated intervals,
including itself, overlap a given interval; for $N=4,6,8$ the number is
even smaller. Bounding the remaining inner products by their norms gives
\begin{equation}
 M_4\le\frac{9}{\ell}\|G_i\|_2^2
      \le576c^2M_2.
 \label{eq:uniform-M4-bound}
\end{equation}
The constants depend on local spin and support size, independently of $N$.

Separate the spectral measure into $|\omega|\le\nu$ and
$|\omega|>\nu$. The low-frequency contribution to $M_2$ is at most
$\nu^2M_0$. Cauchy--Schwarz for the remaining measure therefore gives
\begin{equation}
 C^{\rm tr}_{>\nu}\ge
 \frac{[M_2-\nu^2M_0]_+^2}{2M_4},
 \label{eq:window-moment-inequality}
\end{equation}
where $[x]_+=\max(x,0)$ and the factor two accounts for the two frequency
signs. Taking $\nu=\sqrt{3c/4}$ and using the preceding inequalities,
\begin{align}
 C^{\rm tr}_{>\sqrt{3c/4}}
 &\ge\frac{M_2}{4608c^2}\notag\\
 &\ge\frac{b^2(4c-3)(86c-27)}{3732480\ell}>0
 \quad(s\ge1,\ b\ne0).
 \label{eq:general-window-bound}
\end{align}
For the pentamer quartet, this becomes
\begin{equation}
 C^{\rm tr}_{>\sqrt{45}/4}\ge\frac{197}{1036800}b^2>0
 \qquad(N\ge4\text{ even}).
 \label{eq:quartet-window-bound}
\end{equation}
The bound holds for every real $u$, with the effective uniform-trace
limit taken before increasing $N$. It prevents all of the optimized
weight from vanishing or moving towards zero frequency in this limit.
At fixed positive $\tau$, Eq.~\eqref{eq:thermal-criterion} establishes
nonvanishing weight for each finite chain; a size-uniform bound at that
temperature requires a separate argument.

Full-chain matrices at $s=1,3/2$ and $N=4,6$ verify the zeroth-, second-,
and fourth-moment inequalities. Complete four-site diagonalization also
checks Eq.~\eqref{eq:window-moment-inequality} from actual transition
weights. Minimizing the quartet window weight gives $0.350008b^2$,
about 1842 times the lower bound. The bound consequently establishes
size-uniform nonvanishing weight while allowing substantially larger
model-specific intensities.

\section{Extension to independent cell-local controls}
\label{app:local-controls}

Cell-local exchange changes enlarge the adjustable dipole space while
leaving the leading composite vector fixed. To show this, add
$\delta H_{\rm loc}=\beta\sum_i W_i$, where each bounded Hermitian $W_i$
is an SU(2) scalar confined to cell $i$. The reference ground multiplet is
an isolated single spin-$s$ irrep. Schur's lemma gives
$P_iW_iP_i=w_iP_i$, a constant rather than an effective spin interaction.
The first-order derivative of the local vector probe is again a vector.
Because rank one occurs once in the operator space of a single spin-$s$
irrep, its projection must be a multiple of $T_i^z$. At $q=\pi$ the
probe is Hermitian, so this multiple is real. Changes in the projected
end-spin coefficients enter $\beta V$ only at order $\beta^2$;
likewise, the change of $b$ affects the physical first-order source only
at order $\beta^2$. The leading Hamiltonian is consequently still the
same uniform $h$, and all such controls are contained in the real
one-spin space of Eq.~\eqref{eq:local-control-probe}. Cell reflection
symmetry need not be preserved by $W_i$.
The construction concerns bounded first-order perturbations about the
uniform reference root, with its isolated irrep unchanged.

Let $U$ translate by one effective site. For eigenstates of $h$, define
the positive semidefinite sesquilinear form
\begin{equation}
 (A,B)_{\tau,\nu}=\sum_{\epsilon_f-\epsilon_i>\nu}
 p_i A_{fi}^*B_{fi},\qquad \nu\ge0,
 \label{eq:window-inner-product}
\end{equation}
where $p_i$ are Gibbs weights or $1/d^N$. All degeneracies are included
as complete blocks. Since $[U,h]=0$, conjugation by $U$ is an isometry
of this form. Distinct operator momenta are orthogonal even when the
form has a null space. The $\pi$-momentum projection is
\begin{equation}
 \Pi_\pi A=\frac1N\sum_{m=0}^{N-1}(-1)^mU^m A U^{-m}.
\end{equation}
For even $N$, it obeys
$\Pi_\pi R_{\{u_i\}}=R_{\bar u}$. Pythagoras in
Eq.~\eqref{eq:window-inner-product} gives
Eq.~\eqref{eq:local-control-projection}, and hence
\begin{equation}
 \min_{\{u_i\}\in\mathbb R^N}F(R_{\{u_i\}})
 =\min_{u\in\mathbb R}F(R_u).
 \label{eq:local-control-optimum}
\end{equation}
The equality uses the availability of uniform controls in the enlarged
space. Any more restricted microscopic control set still satisfies the
same lower bound. The nonzero derivative $A'$ makes the uniform direction
accessible in the pentamer example. This argument applies at every
positive scaled temperature and to the uniform-trace window bound;
it does not require a nondegenerate spectrum.

The second moment provides an additional explicit check. On bond
$(i,i+1)$ the dipole commutator coefficient becomes
$(u_i+u_{i+1})/2$. Repeating the support-orthogonality calculation of
Appendix~\ref{app:general-moment} gives
\begin{align}
 M_2(\{u_i\})={}&M_{2,\min}+\frac{(d^2-1)^2}{18\ell N}\notag\\
 &\times\sum_i\left[\frac{u_i+u_{i+1}}2+
              \frac{b(d^2-4)}{12}\right]^2.
 \label{eq:nonuniform-moment}
\end{align}
An alternating component of $u_i$ adds a conserved total $T^z$ to the
probe and leaves this moment unchanged, although it changes $M_0$.
For nonuniform controls, the finite-frequency lower bound therefore
follows from the operator-momentum decomposition rather than the
uniform-control $M_2/M_0$ inequality.

For all ten pairwise exchange directions inside the isolated pentamer,
the derivative $-P(W_iGB+BGW_i)P$, with
$G=Q(H_{\rm cell}-E_0)^{-1}Q$, has a non-dipole Frobenius norm below
$2.4\times10^{-15}$. Central differences at step $5\times10^{-5}$
agree within $2.8\times10^{-10}$. Full four-site spectra for $s=1,3/2$
at $\tau=0.2,1.6$ and uniform trace verify
Eq.~\eqref{eq:local-control-optimum} for $\nu=0$ and
$\sqrt{3s(s+1)/4}$, with maximum difference $1.7\times10^{-16}$ at
$b=1$. Full-chain matrices at $N=4,6$ verify
Eq.~\eqref{eq:nonuniform-moment} within $1.5\times10^{-14}$.

\section{Microscopic tests and finite-chain controls}
\label{app:ground-controls}

Zero-temperature spectra test the state dependence of exchange
cancellation and separate the composite response from endpoint effects.
The comparisons below combine microscopic exact diagonalization with
effective-chain calculations at matched local roots.

\subsection{Different powers on the same periodic geometry}

Local diagonalization found simple roots at each of 16 sampled values
of $Q/\pi$ from $0.55$ to $1$, using $q=(Q+4\pi)/5$ for the quartet
and $q=(Q+2\pi)/5$ for the doublet [Fig.~\ref{fig:rank-controls}(a)].
All sampled local gaps were positive, with minimum $0.04263$.
The quartet coefficient $b$ was nonzero throughout this sample,
whereas the doublet quadrupole vanished identically.
These local roots provide candidate interference settings; a nonzero
many-body residual must be established separately at each setting.

The microscopic comparison uses the same $N=3$, $L=15$ periodic
geometry for both local spins. Table~\ref{tab:rank-roots} specifies simple
roots at different physical momenta but the same effective momentum
$Q=2\pi/3$. Both roots belong to the Hamiltonian in
Eq.~\eqref{eq:rank-ham}.
The odd periodic geometry is frustrated, so this comparison tests the
finite-system mechanism rather than an infinite unfrustrated chain.

\begin{table}[!t]
\caption{Local roots and predictions for the microscopic three-cell comparison.
The local gap is evaluated at the root. $C_0=\|k\|^2$ and $\Cmin$ are
effective predictions for the coefficient of $\beta^2$, not fits to the
microscopic weights.}
\label{tab:rank-roots}
\begin{ruledtabular}
\begin{tabular}{lrr}
 & Quartet & Doublet\\
$s$ & $3/2$ & $1/2$\\
$q/\pi$ & $14/15$ & $8/15$\\
$\alpha_*$ & $-0.38103454$ & $1.45352917$\\
$\Dloc$ & $0.10336638$ & $1.10511693$\\
$\kappa_{\mathrm{opt}}$ & $-9.45913710$ & $-0.17229706$\\
$C_0$ & $0.54511625$ & $0.000474008$\\
$\Cmin$ & $0.14202789$ & $<10^{-30}$\\
\end{tabular}
\end{ruledtabular}
\end{table}

We diagonalized the microscopic $S^z=1/2$ sector of dimension 6435 at
$g=0.005,0.01,0.02,0.04$, both at the fixed local root and with its
predicted linear retuning. There are two degenerate ground states in this
sector; their equal average excludes the entire ground eigenspace from the
inelastic weight. All ten quartet or one doublet low-band excited states
in this sector were included, with additional states above the cutoff used
to check band completeness. The largest eigen-residual was below
$3.1\times10^{-13}$.

The retuned quartet weight remained quadratic at weak coupling
[Fig.~\ref{fig:rank-evidence}(a)]. At the smallest $g$, its coefficient
$\Wlow/\beta^2=0.14091678$ differs from the independent prediction
$0.14202789$ by $0.7823\%$. A log--log fit over the three smallest couplings
gave a slope $1.9869$. For the doublet the corresponding slope was $4.0167$
[Fig.~\ref{fig:rank-evidence}(b)]. Its $\Wlow/\beta^4$ varied from
$1.0694\times10^{-4}$ to $1.1290\times10^{-4}$ across all four couplings.
These finite-interval slopes test perturbative orders along the specified
retuning trajectories. The doublet result supports
Eq.~\eqref{eq:rank-doublet}; further retuning may suppress its weight
below fourth order.

The cancellation was also checked at the operator level. In the full
eight-dimensional retained doublet space, the relative norm of
$K_q+\kappa_{\mathrm{opt}}A'D_Q$ was $4.32\times10^{-15}$.
The analogous quartet operator did not vanish in its full 64-dimensional
space. A restart with a different random vector, larger Krylov space, and
tighter tolerance reproduced the two weakest doublet weights to relative
differences below $1.3\times10^{-9}$.

\begin{figure*}[t]
\centering
\includegraphics[width=0.98\textwidth]{rank_compensation.pdf}
\caption{Retuned-weight comparison and validation limits.
(a,b) Absolute microscopic low-band weights on the same three-cell periodic
geometry, at the parameters in Table~\ref{tab:rank-roots}. Dashed curves keep
$\alpha=\alpha_*$; solid curves use $\alpha_*+\kappa_{\mathrm{opt}}\beta$.
The vertical ranges differ, but no independent intensity normalization is
applied. Points are calculated values; segments only connect them.
(c) Leading quartet minimum $\Cmin$ at the exact $q=\pi$ root for even
effective chains. ED denotes exact diagonalization, OBC/PBC denote open/periodic
boundaries. The accepted DMRG point is an effective $N=20$ open chain;
the $N=40$ cross failed the preset energy-convergence criterion and is not
used as converged evidence. No thermodynamic fit is drawn.
(d) Microscopic $L=20$ open-chain comparison of the retuned leading probe
with the line-weight vector and integrated weight. The errors at larger
$g$ reveal failure of the leading theory. All panels use calculated data
without smoothing; logarithmic axes in (a,b,d) change only the display.}
\label{fig:rank-evidence}
\end{figure*}

\subsection{Even periodic chains and the open-end contribution}

The six-cell effective ring uses the same local roots and effective
momentum $Q=2\pi/3$ as Table~\ref{tab:rank-roots}.
The quartet effective minimum was $0.06641159$, while the doublet minimum
was below $4\times10^{-32}$. The quartet's optimum changed to
$\kappa=-5.56848$, reflecting its dependence on the many-body ground state;
the doublet's operator-level optimum remained $-0.17229706$.
Thus the contrast also occurs without odd-ring frustration in the effective
model. This is not a microscopic thirty-spin calculation.
At the same quartet root, the accepted $N=12$ periodic MPS gives
$\Cmin=0.08315075$ and $\kappa_{\mathrm{opt}}=-5.67393$
[Fig.~\ref{fig:rank-controls}(d)]. Increasing its bond dimension from
384 to 512 changes this weight by only $6.15\times10^{-7}$ relatively.
Its scalar contribution is removed to within $1.4\times10^{-17}$ in
the residual norm squared, again identifying a composite rather than
endpoint source. These size-dependent coefficients are not extrapolated.

For periodic effective chains at the exact quartet root in
Eq.~\eqref{eq:rank-exact-root}, the sizes $N=6,8,10$ gave
$\Cmin=0.000443300,0.000649072,0.000751090$, respectively.
However, $N=4$ gave zero to numerical precision
[Fig.~\ref{fig:rank-evidence}(c)]. This exception rules out attributing an
unavoidable nonzero residual to the local spin alone.

Periodic DMRG extends this test to $N=12$, where bond dimension 512 gives
$\Cmin=0.00080811794$. Relative to bond dimension 384, the coefficient
changes by $3.52\times10^{-5}$ and the energy per cell by
$2.97\times10^{-9}$, passing the original convergence thresholds.
The energy variance per cell is $6.73\times10^{-10}$, and the bond-energy
spread is $1.41\times10^{-8}$. These periodic results exclude an endpoint
origin for this finite-chain residual without establishing a thermodynamic
lower bound.
The $N=16$ calculation at bond dimension 768 had a stable weight but an
energy change per cell of $1.304\times10^{-8}$, above the preset
$10^{-8}$ gate. It is marked unaccepted in Fig.~\ref{fig:rank-controls}(c)
and is not used to extend the convergence claim.

To distinguish bulk-link physics from endpoints, we separated $K_q$ into the
one-spin and quadrupole-times-spin pieces of Eq.~\eqref{eq:rank-link} before
projecting each perpendicular to $d$. For all periodic cases, the perpendicular
one-spin norm squared was below $2\times10^{-30}$. The nonzero residuals are
therefore due to the composite source. Open chains do not share this property.
At $N=10$ OBC, the one-spin, composite, and interference contributions to
$\Cmin$ were $0.000286554$, $0.001244067$, and $-0.000513469$.
Their sum is $0.001017151$. Calling this whole open-chain residual a pure
composite response would be incorrect.

Retuning also changed which excitations dominated the remaining effective
weight. For $N=10$ PBC, the lowest-line share decreased from $0.94007$ at the
fixed root to $1.04\times10^{-5}$ at the integrated-weight optimum. The mean
energy of the residual, in units of $\beta$, rose from $0.09456$ to $0.85896$.
These numbers describe an effective finite chain. They do not imply a bulk
continuum, and the integrated optimum is not exactly the value that cancels
the lowest line alone.

An independent MPS calculation reproduced the $N=10$ open-chain ED minimum
to relative error $2.4\times10^{-6}$ at bond dimension 128. Refinement to
bond dimension 256 gave an accepted $N=20$ minimum $0.00100937976$.
The $N=40$ weight appeared stable but its energy change per cell,
$1.0814\times10^{-8}$, remained above the preset $10^{-8}$ threshold.
It is reported as unaccepted, not used to infer a nonzero thermodynamic limit.

\subsection{Microscopic validity at the exact quartet root}

\begin{figure*}[t]
\centering
\includegraphics[width=0.98\textwidth]{rank_controls.pdf}
\caption{Local mechanism and periodic-chain controls.
(a) Simple local roots at 16 common effective momenta, with
$q=(Q+4\pi)/5$ for quartets and $q=(Q+2\pi)/5$ for doublets. Segments connect
calculated roots; this is not a many-body phase diagram.
(b) Exact channel coefficient $r_J$ in Eq.~\eqref{eq:bond-channel}, not a
spectral intensity. Different quartet channels cannot all be compensated by
the same scalar shift at the link-operator level.
(c) Quartet $\Cmin$ at $q=\pi$ on even periodic effective chains.
The $N=12$ DMRG point passed the original convergence tests; the $N=16$
cross did not and is not converged evidence. No bulk fit is drawn.
(d) Effective chains at the same roots and $Q=2\pi/3$ as
Table~\ref{tab:rank-roots}; the quartet optimum is recomputed for each size.
The doublet dashed line is the first-order operator-cancellation result,
with ED dots at $N=3,6$, not a microscopic finite-$\beta$ zero.
Panels (c,d) retain the normalization $1/\sqrt{5N}$ but use different physical
momenta and local roots; they are not one common size series.
All weights are computed without smoothing or reflection averaging.}
\label{fig:rank-controls}
\end{figure*}

Four open pentamers provide a second microscopic test at $q=\pi$, without
the frustration of the three-cell ring. The effective prediction for their
retuned coefficient is $\Cmin=0.00067813935$.
At $g=0.000625$, microscopic diagonalization gave $0.00067904306$, an
integrated-weight error of $0.1333\%$. The relative Euclidean error of the
sorted low-band line-weight vector was $1.2254\%$.
Doubling $g$ to $0.00125$ approximately doubled that line error to
$2.4625\%$; an independent solver restart at the smaller $g$ agreed to
relative line error $3.56\times10^{-9}$.

The agreement is restricted to weak coupling. At $g=0.02$ and $0.04$, the
line-vector errors were $46\%$ and $109\%$, and the integrated-weight errors
were $14\%$ and $50\%$ [Fig.~\ref{fig:rank-evidence}(d)]. These results
prevent extending the leading composite-probe interpretation quantitatively
to the much stronger couplings of the original full $L=40$ spectra.

\begin{figure*}[t]
\centering
\includegraphics[width=0.49\textwidth]{fig6_fullq_positive.pdf}\hfill
\includegraphics[width=0.49\textwidth]{fig6_fullq_negative.pdf}
\caption{Complete connected spectra for $L=40$ open chains at finite coupling.
$J_2/J_1=0.5$ (left), $-0.5$ (right), and $J_3/J_1=0.5$.
The triangular time window at $TJ_1=60$ has width $0.09277J_1$;
the production time step and bond dimension are $\Delta tJ_1=0.05$ and 192.
Both use the same display coordinate
$u(s)=\operatorname{asinh}(s/0.02)/\operatorname{asinh}(6/0.02)$ for
$s=J_1S_T^{zz}\in[0,6]$. This monotonic display does not change the spectral
arrays or their resolution. No Gaussian smoothing, momentum interpolation,
or reflection averaging is applied. The 41 momenta were independently
calculated. These finite-coupling maps are contextual, not quantitative
evidence for the weak-coupling retuned residual.}
\label{fig:rank-original}
\end{figure*}

\section{Local root and numerical protocols}

For the quartet root at $q=\pi$, the highest-weight $M=3/2$ state has
single-down-spin amplitudes $(u,v,w,v,u)$, with $w=-2(u+v)$ from the
highest-weight condition. The zero of the projected alternating probe gives
$v^2=1/8$ after normalization. The relevant eigenvector ratio is
$x=u/v=(\sqrt7-2)/3$, with
$\alpha=-x(3+2x)/(1-x^2)$, giving Eq.~\eqref{eq:rank-exact-root}.
Direct diagonalization of the 32-dimensional pentamer Hamiltonian
independently supplies the local gaps and resolvents.

Microscopic Hamiltonians were built in fixed-$S^z$ bases. The three-cell
comparison used a real symmetric sparse eigensolver with tolerance
$10^{-13}$, at least 50 Krylov vectors, and four extra eigenpairs above the
expected retained band. The numerical acceptance tests required two ground
states in the chosen sector, a complete retained band, and eigen-residuals
below $10^{-10}$. The anticipated small-coupling slope intervals were
$[1.8,2.2]$ and $[3.7,4.3]$; the quartet's weakest-coupling coefficient was
required to agree within $5\%$. These are deterministic validation tests,
not uncertainty intervals. All passed. The microscopic tests use fixed
roots and retuning coefficients derived from the effective probe.
The restart check used tolerance $10^{-14}$, 120 Krylov vectors, a different
seed, and six extra eigenpairs above the retained count at $g=0.005,0.01$.
This solver-sensitivity check uses the same Hamiltonian assembly.

For the four-cell quartet test, the $S^z=0$ microscopic basis has 184756
states. We obtained 45 eigenpairs, including all 44 retained-sector states
(the ground state and 43 excited states), with tolerance $10^{-13}$.
The two smallest couplings used 140 Krylov vectors; a repeat at
$g=0.000625$ used 180 and a different random starting vector. The line-vector
error compares weights in increasing-energy order, including a zero
ground-state entry, after division by $\beta^2$. It is not a test of a
broadened line shape. The strongest-coupling points were diagnostic tests,
not part of the weak-domain acceptance claim.

Effective-chain ED used the $S^z=0$ sector, whose dimension reaches 116304
for ten spin-$3/2$ sites. The norm and overlap in
Eq.~\eqref{eq:rank-minimum} were evaluated directly on the ground-state
sources, without summing an incomplete set of excited states. The lowest
line and the energy moment were checked separately. The periodic scalar
source and the sum of the open-chain decomposition supplied internal
operator checks.

The longer open effective chains were computed by DMRG~\cite{white1992,schollwoeck2011}
using ITensors/ITensorMPS~\cite{fishman2022}
with conserved $S^z$, explicit spin-$3/2$ local operators, and MPO contractions
of $D^\dagger D$, $K^\dagger K$, and $D^\dagger K$. Thus there is no additional
truncation from constructing an applied source MPS. Refinement used bond
dimensions $128,192,256$, with 24, 16, and 16 sweeps, cutoff $10^{-13}$,
and eigensolver tolerance $10^{-13}$. Larger-bond runs started from the
previous state without decreasing its bond dimension. Acceptance required
relative change of $\Cmin$ below $1\%$, energy change per cell below
$10^{-8}$, energy variance per cell below $10^{-8}$, and maximum local
$|\langle T_n^z\rangle|<10^{-4}$. The final $N=20$ variance was
$1.21\times10^{-9}$ and its maximum local moment was $1.17\times10^{-9}$.
The earlier unconverged runs and the failed $N=40$ refinement are retained
with their original criteria.

Periodic DMRG used the same operators with the closing bond included in the
MPO. The $N=6$ source norm was cross-checked against ED. Bond dimensions
$128,256,384$ were followed by 512 for $N=12$, with cutoff $10^{-13}$ and
16 warm-start sweeps at each refinement. In addition to the preceding
energy and weight criteria, acceptance required
$|\langle\bm T_{\mathrm{tot}}^2\rangle|<10^{-4}$, bond-energy spread below
$10^{-5}$, and scalar-removal and energy-reconstruction identities within
$10^{-10}$. The first $N=12$ batch failed the energy-change test at 384;
the subsequent 512 calculation passed without changing that threshold.
All sizes here count effective spin-$3/2$ sites, not microscopic spins.
The bounded refinement also tested $N=16$ at dimensions $256,512,768$;
the last run failed the unchanged energy-change threshold and no larger
periodic calculation was started.

For the $Q=2\pi/3$ quartet check, the same periodic ground states can be
reused: changing the local root multiplies the leading uniform Heisenberg
Hamiltonian by the positive scalar $c_1(\alpha)^2$ without changing its
eigenvectors. The probe coefficients were recomputed at the new root,
including their complex Fourier phases. Connected norms use
$\langle D^\dagger K\rangle-\langle D\rangle^*\langle K\rangle$,
and the exchange retuning remains real. The six-cell MPO result reproduced
ED to relative error $4.5\times10^{-15}$. The twelve-cell weight passed the
$1\%$ bond-dimension-change test, and the imaginary part of the normalized
source overlap was below $10^{-8}$. The calculation measures the
effective-model response on twelve cells.

\section{Thermal traces and exact local checks}
\label{app:exact-local-check}

The four- and six-cell thermal effective calculations included every
$S^z$ sector, with total dimensions $4^N$ or $2^N$. For six quartet
cells the largest block had dimension 580. Real and imaginary operator
components were transformed into the complete energy bases, and all
transitions with $|\epsilon_f-\epsilon_i|<10^{-9}$ were treated as
equal-energy transitions. Checks included Hermiticity, trace sum rules,
eigen-residuals below $10^{-11}$, detailed balance, and agreement between
direct residual norms and the thermal least-squares expression. A separate
256-dimensional tensor-product assembly reproduced the four-cell result.
Random unitary changes within each degenerate eigenspace left the thermal
weights invariant. Both floating-point assemblies used the same dense
diagonalization library.

\paragraph*{Exact microscopic prefactors.}
At the quartet root, an unnormalized highest-weight cell state in the
one-down-spin basis, ordered by the position of that spin, is
\begin{equation}
 v=(1,\ 2+\sqrt7,\ -6-2\sqrt7,\ 2+\sqrt7,\ 1)^{\mathsf T}.
\end{equation}
Its squared norm is $88+32\sqrt7$ and its energy is
$E_5=-(5+4\sqrt7)/12$. Total-spin lowering constructs the quartet
projector $P_{\rm cell}$. We evaluated the reduced resolvent
$G=(H_{\rm cell}-E_5+P_{\rm cell})^{-1}-P_{\rm cell}$ separately in all
six magnetization blocks, using exact arithmetic in $\mathbb Q(\sqrt7)$.
An exact $LDL^{\mathsf T}$ factorization of
$H_{\rm cell}-E_5-\tfrac18(1-P_{\rm cell})+P_{\rm cell}$
has strictly positive diagonal pivots in every block. This certifies
that the quartet is the ground multiplet and that $\Dloc>1/8$.
The projected tensor in Eq.~\eqref{eq:rank-local-tensor} is
$s_0\delta_{z\mu}I+t_0Q^{z\mu}$, with
\begin{align}
 s_0&=\frac{161590\sqrt7-361001}{180792},\notag\\
 t_0&=\frac{169277-70390\sqrt7}{45198}.
\end{align}
Both endpoint choices and all $z,+,-$ components satisfy this identity
exactly. Multiplication by $-c_1$ gives the coefficients in
Eq.~\eqref{eq:rank-link},
\begin{align}
 a&=\frac{4500505-2595746\sqrt7}{19525536},\notag\\
 b&=\frac{1082642\sqrt7-2261005}{4881384}>0.
 \label{eq:exact-local-coefficients}
\end{align}
The elementary bounds
$2.645<\sqrt7<2.646$ certify $b>20086103/162712800>0$.
Comparison with the separately assembled floating-point two-cell
resolvent gives an absolute coefficient difference below
$3.0\times10^{-16}$ for the same perturbative formula.

The other exact checks used rational arithmetic in the unnormalized local basis
$|k\rangle=(T^-)^k|s,s\rangle$, where
$T^+|k\rangle=k(2s-k+1)|k-1\rangle$ and
$\langle k|k\rangle=k!(2s)!/(2s-k)!$.
This gives exact zero tests for the partial traces in
Eq.~\eqref{eq:thermal-bound}, with no numerical threshold.
For the four-cell activation calculation, the two sublattice spins were
coupled from $j_A=j_B=3$. In an unnormalized product basis, a highest-weight
total-spin-$J$ state has coefficients $c_k$ at sublattice lowering degrees
$(k,6-J-k)$, obeying
\begin{equation}
 \frac{c_{k+1}}{c_k}
 =-\frac{(6-J-k)(J+1+k)}{(k+1)(6-k)}.
\end{equation}
This recurrence, followed by total-spin lowering, constructs the ground
singlet and first triplet without a floating-point eigensolver.
Acting with $Y+3\widetilde D$ annihilates the singlet. On the triplet,
the residual separates into final eigenvalues $-9$ and $-4$ of $h$;
projection by $\bm A^2$ resolves the two $-4$ channels. Their exact norms
give Eq.~\eqref{eq:thermal-activation}. The leading coefficient agrees
with the separate full-space thermal calculation to relative error
$3\times10^{-15}$.

The ten-spin microscopic thermal calculation used all 11 $S^z$ blocks,
whose largest dimension is 252. The full partition function was retained.
Both initial and final states in $W_+$ were restricted to
$E-E_0<\Dloc/2$; there were 16 retained quartet or 4 retained doublet
states, as required. Eigen-residuals were below $3.4\times10^{-15}$,
and omitted-band thermal populations were below $1.1\times10^{-57}$
for the tested parameters. The temperature was
$T=\tau\beta c_1(\alpha_*)^2$, including on the retuned trajectory.

For $g=0.005,0.01,0.02$, the retuned quartet weights have logarithmic
slopes $2.07089$ and $2.02549$ at $\tau=0.4$ and $1.6$;
the doublet slopes are $4.00185$ and $4.00189$.
At the weakest coupling, the quartet coefficients differ from the
effective predictions by $1.75\%$ and $1.18\%$, respectively.
At $\tau=0.4,g=0.02$, the deviation reaches $12.25\%$,
showing the finite-coupling correction to the leading response.

For the finite momentum interval we formed the low-band matrix elements
of each microscopic $S_j^z$. At each quadrature point the complex Fourier
amplitudes were combined before taking their absolute squares; this
avoids subtracting nearly equal contributions to a tiny central intensity.
Uniform-window integrations with 16 and 32 Gauss--Legendre nodes agreed
within $10^{-7}$ relatively. Zero-width values reproduced the separate
microscopic calculation within $1.2\times10^{-10}$. This is integration
over a specified momentum acceptance, not Gaussian broadening or an
image-processing operation.

For twenty microscopic spins, the sparse $S^z=0$ Hamiltonian has dimension
184756. We computed 48 lowest eigenpairs and verified that precisely 44
lie below $E_0+\Dloc/2$. Within this low band, diagonalization of
$(S^+V)^\dagger(S^+V)$ resolves total spin. Rediagonalizing the Hamiltonian
within each spin sector and applying normalized total-spin raising gives
all positive-$S^z$ states; spin reversal supplies the negative sectors.
The multiplicities of $J=0,\ldots,6$ are $4,9,11,10,6,3,1$,
giving 256 states. Every raised sector was checked against an independently
assembled microscopic Hamiltonian. Residuals were below $10^{-12}$ and
orthogonality errors below $10^{-11}$. The same procedure at ten spins
reproduced the full-spectrum thermal weight to relative error below
$10^{-7}$.
At the weakest coupling, an independent restart with a different random
vector, a larger Krylov space (220 instead of 180), and tighter tolerance
changed the total, ground-initial, and excited-initial weights by at most
$8.0\times10^{-8}$ relatively across the tested temperatures and both
frequency definitions, using the same microscopic Hamiltonian.

The omitted states have a rigorous thermal bound. Relative to uncoupled
cells at $\alpha_*$, the perturbation norm is at most
$v=(3/4)[N|\beta|+2N|\alpha-\alpha_*|]$.
The verified condition $4v<\Dloc$ isolates the retained band and places
it below the chosen cutoff. Weyl bounds put every omitted state at least
$\Dloc-2v$ above the true ground energy. With the retained partition
function shifted by that energy,
\begin{equation}
 \frac{Z_{\rm out}}{Z_{\rm low}}
 \le \frac{(2^{5N}-4^N)e^{-(\Dloc-2v)/T}}{Z_{\rm low}}.
\end{equation}
The logarithm of this bound is below $-2167$ in all tested thermal cases;
the tail was bounded, not set to zero by numerical underflow.
The unwindowed positive-frequency cutoff was $10^{-10}$ in $J_1$ units.
The unwindowed sums include quasielastic transitions and were checked
with cutoffs $10^{-11}$ and $10^{-9}$ in $J_1$ units.
The regular response uses $\omega>\eta J_{\rm eff}$ with $\eta=0.05$;
changing $\eta$ from $0.02$ to $0.1$ leaves the computed window weights
unchanged. At the weakest coupling, the windowed coefficients agree
with the effective predictions to $0.161\%$, $0.229\%$, and $0.042\%$
at $\tau=0.2,0.4,1.6$. Both frequency definitions are retained in the
numerical tables.

\section{Finite-coupling spectral context}

The $L=40$ maps in Fig.~\ref{fig:rank-original} show the redistribution
of spectral weight between the doublet and quartet regimes at
$\beta=0.5$. They have
$\beta/\Dloc\simeq1.46$ and $3.96$, outside the validated perturbative
range, and finite-time resolution broader than the separately calculated
open-chain gaps. They neither establish the weak-coupling powers in
Fig.~\ref{fig:rank-evidence} nor resolve the tiny retuned weights; they
are retained as context, not validation of the composite-probe expansion.


\bibliography{references,rank_references}
\end{document}
